% ==================================================================
% id:       seq-013
% topic:    סדרות
% subtopic: גבולות חלקיים
% solution: full
% ==================================================================
\textbf{בכל אחד מהסעיפים הבאים, מצאו את כל הגבולות החלקיים במובן הרחב של הסדרה הנתונה:}
\bE\bMs{2}
 \item $\ds{a_n=\cos\lr{\frac{\pi n}{6}}}$
\hint{חשבו כמה איברים ראשונים ושימו לב שהסדרה מחזורית. מה אורך המחזור?}
\begin{solbox}
נחשב תחילה כמה איברים בגבולות המחזור הראשון:
\[
\begin{aligned}
a_1  &= \cos\left(\frac{\pi}{6}\right)  = \frac{\sqrt{3}}{2},\quad
a_2  = \cos\left(\frac{\pi}{3}\right)  = \frac12,\\
a_3  &= \cos\left(\frac{\pi}{2}\right)  = 0,\quad
a_4  = \cos\left(\frac{2\pi}{3}\right) = -\frac12,\\
a_5  &= \cos\left(\frac{5\pi}{6}\right) = -\frac{\sqrt{3}}{2},\quad
a_6  = \cos(\pi)                      = -1,\\
a_7  &= \cos\left(\frac{7\pi}{6}\right) = -\frac{\sqrt{3}}{2},\quad
a_8  = \cos\left(\frac{4\pi}{3}\right) = -\frac12,\\
a_9  &= \cos\left(\frac{3\pi}{2}\right) = 0,\quad
a_{10} = \cos\left(\frac{5\pi}{3}\right)  = \frac12,\\
a_{11} &= \cos\left(\frac{11\pi}{6}\right) = \frac{\sqrt{3}}{2},\quad
a_{12} = \cos(2\pi)                     = 1.
\end{aligned}
\]
בגלל שאיברי הסדרה מחזוריים
$$a_{n+12}=\cos\left(
\frac{\pi(n+12)}{6}\right)=\cos\left(
\frac{\pi n}{6}+2\pi\right)=\cos\left(
\frac{\pi n}{6}\right)=a_n$$
הגבולות החלקיים יתקבלו מתתי-הסדרות
$a_{12n},a_{12n+1},\ldots,a_{12n+11}$
ולכן יהיו
 $$.\left\{-1,-\frac{\sqrt{3}}{2},-\frac{1}{2},0,\frac{1}{2},\frac{\sqrt{3}}{2},1\right\}$$
\end{solbox}
 \item $\ds{a_n=\frac{(-1)^n}{n}}$
 \begin{solbox}
זוהי סדרה מתכנסת (מדוע?), לכן גבולה החלקי היחיד הוא הגבול שלה שהוא
אפס.
\end{solbox}

 \item $\ds{a_n=\lr{1+(-1)^n}\cdot n}$
\hint{פצלו לפי הזוגיות של $n$.}
\begin{solbox}
נציג
$$.a_n=\case{cc}{2n &\text{\R{ זוגי }}n\\0 & \text{\R{ אי-זוגי }}n}$$
מכיוון ש-
$\displaystyle\lim_{n\to\infty}a_{2n}=\infty$
ו-
$\displaystyle\lim_{n\to \infty}a_{2n+1}=0$,
נסיק
שגבולותיה החלקיים במובן הרחב הם
$0,\infty$
וכי
 לא קיים הגבול גם במובן הרחב.
\end{solbox}
 \item $\ds{a_n=(-1)^n\sqrt[n!]{n^{100}}+\frac1{n^2}\cos\lr{\frac{n}2}}$
\hint{טפלו בכל מחובר בנפרד.}
\begin{solbox}
הסדרה
$\frac1{n^2}\cos\lr{\frac{n}2}$
שואפת ל-
$0$
כמכפלה של סדרה אפסה בחסומה, לכן גבולה החלקי היחיד הוא פשוט הגבול שלה והוא, כאמור,
$0$.
בנוסף, מתקיים
$$,1\leq n^{\frac{100}{n!}}<n^{\frac{100}{n}}
=(\sqrt[n]{n})^{100}\xrightarrow[n\to\infty]{}1$$
לכן, מכלל סנדוויץ', נסיק כי גבולה של
$n^{\frac{100}{n!}}$
הוא
$1$.
נסכם כי בעזרת אריתמטיקה מתקבל
$a_{2n}\xrightarrow[n\to\infty]{}1$
וגם
$a_{2n+1}\xrightarrow[n\to\infty]{}-1$
ולכן לסדרה הנתונה שני גבולות חלקיים בלבד והם
$1$
ו-
$-1$.
\end{solbox}
 \item $\ds{a_n=(-1)^n\cdot2^{\frac1{2n-1}}\cdot\sqrt[n!]{n^{2025}}+\frac1{n}\sin\lr{e^n}}$
\hint{השתמשו בסעיף הקודם, ובדקו את הגבול של כל גורם בנפרד.}
\begin{solbox}
הסדרה
$\frac1{n}\sin\lr{e^n}$
שואפת ל-
$0$
כמכפלה של סדרה אפסה בחסומה, לכן גבולה החלקי היחיד הוא פשוט הגבול שלה והוא, כאמור,
$0$.
בסעיף הקודם ראינו כי
$\displaystyle\lim_{n\to\infty}\sqrt[n!]{n^{100}}=1$,
ובצורה זהה גם
$\displaystyle\lim_{n\to\infty}\sqrt[n!]{n^{2025}}=1$.
מכיוון שהסדרה
$2^{\frac{1}{2n-1}}$
שואפת ל-
$1$,
נסיק כי
$\pm1$
הם גבולותיה החלקיים היחידים של הסדרה הנתונה, בדומה לסעיף הקודם.
\end{solbox}
 \item
 $\ds{a_n=\frac{\left((-1)^n-1\right)n}{\sqrt{n+1}}+
\frac{(-1)^{n}n^3-4}{2n^3}}
$
\hint{פצלו לפי הזוגיות של $n$ ובדקו מה קורה למחובר הראשון בכל מקרה.}
\begin{solbox}
כאן מתקיים
$$a_{2n}=\frac{(2n)^3-4}{2(2n)^3}=\frac{8n^3-4}{16n^3}\rightarrow \frac{8}{16}=\frac{1}{2}$$
וגם
$$a_{2n+1}=\frac{-2(2n+1)}{\sqrt{2n+2}}+\frac{-(2n+1)^3-4}{2(2n+1)^3}\rightarrow -\infty$$
כאשר הגבול האחרון נובע מאירתמטיקה במובן הרחב, כי המחובר הראשון שואף למינוס אינסוף והמחובר השני שואף לגבול
$-1/2$.

לכן, הגבולות החלקיים של הסדרה הם
$\frac{1}{2}$
ו-$(-\infty)$.
\end{solbox}
\eMs
\eE
